\documentclass[9pt]{IEEEtran}

\usepackage[english]{babel}
\usepackage{graphicx}
\usepackage{url}
\usepackage{booktabs}
\usepackage{array}
\usepackage{hyperref}
\usepackage[T1]{fontenc}
\usepackage[utf8]{inputenc}
\usepackage{lmodern}
\input{glyphtounicode}
\pdfgentounicode=1

\title{\vspace{0ex}
IEPS: Programming Assignment 3\\
Explainable Retrieval-Augmented Generation}
\author{Dino Džaferagić, Tibor Koderman, Jerneja Krajcar}\vspace{-4.0ex}
\date{\today}

\begin{document}
\maketitle

\section{Introduction}
This assignment extends our PA2 Slovenian information-retrieval system into
an explainable Retrieval-Augmented Generation (RAG) pipeline.  We retain the
successfully evaluated MedOverNet corpus and retrieval configuration, and add
answer generation in two directly comparable modes: \textit{With Context},
where retrieved chunks form explicit evidence, and \textit{Without Context},
where the same local language model receives only the question.

The implementation is intentionally experimental rather than deployed as a
user interface.  For every RAG answer it stores the evidence passages,
original vector ranks, reranked ranks and scores, allowing answer errors to be
related to retrieval failures rather than treated as opaque model behavior.

\section{Implementation Summary}

\subsection{Existing PA2 Evidence Store}
PA3 reuses the PA2 PostgreSQL/pgvector database without schema changes or
re-indexing.  The corpus contains Slovenian health and wellness material from
MedOverNet, including editorial articles and forum threads.  PA2 created
\texttt{page\_segment\_long} chunks: articles are grouped into approximately
250 words with 50-word overlap; forum content is kept close to individual
posts using approximately 220 words with 30-word overlap.

This was the appropriate continuation of PA2: long segments preserve enough
semantic context for a generated response, while PA2 manual evaluation showed
that the 50-character comparison chunks frequently lose meaning at chunk
boundaries.

\subsection{Models and Libraries}
Queries are embedded using \texttt{sentence-transformers/LaBSE}, a
multilingual sentence embedding model whose 768-dimensional output exactly
matches the already populated \texttt{vector(768)} columns.  Candidate
segments are obtained using cosine distance through pgvector and then
reranked with \texttt{BAAI/bge-reranker-v2-m3}, the multilingual
cross-encoder already evaluated in PA2.

Answers are generated through a local Ollama service using
\texttt{gemma3:4b}.  The model has approximately four billion parameters; its
Ollama package is approximately 3.3 GB and its documented multilingual
support makes it a practical local model for Slovenian questions
\cite{ollamagemma}.  A deterministic experimental setting is used with
temperature set to zero.

The implementation is divided into:
\begin{itemize}
    \item \texttt{retriever.py}: LaBSE query embedding, cosine retrieval of 20
    candidate chunks, and reranking to the final top five evidence chunks;
    \item \texttt{prompts.py} and \texttt{generator.py}: separate prompts for
    evidence-grounded and LLM-only answering, and the Ollama API integration;
    \item \texttt{rag\_pipeline.py}: single-question command-line interface;
    \item \texttt{evaluate.py}: batch comparison runner storing full JSON
    traces and manual scoring fields.
\end{itemize}

\subsection{Prompt Design}
In RAG mode, the prompt includes numbered chunks and source URLs.  The model
is instructed to use only that evidence, cite used chunks as \texttt{[1]},
\texttt{[2]}, and explicitly state when the context is insufficient.  This is
particularly important for medical content and out-of-domain questions.

In LLM-only mode, no chunks or URLs are provided.  The model is asked for a
short Slovenian response and told not to invent references.  Keeping the
question and generation model fixed while changing only the evidence enables
a transparent comparison.

\section{System Evaluation Criteria}

\subsection{Data Processing}
No new data preparation is performed in PA3.  The evaluation relies on the
PA2 decision to use logical long chunks with overlap.  Data-processing quality
is therefore assessed through inspectability of the retrieved chunk: a good
chunk should contain an interpretable passage that directly supports an
answer rather than only a matching keyword.

\subsection{Retrieval}
For each query, LaBSE cosine search fetches 20 initial candidates from
\texttt{page\_segment\_long}.  The cross-encoder reranker chooses five
passages supplied to the LLM.  We store the vector rank, cosine distance,
final rank, reranker score, text and URL for every final passage.

Retrieval relevance is scored manually on a three-point scale:
\texttt{0} (not informative), \texttt{1} (partially informative), and
\texttt{2} (directly supports answering the question).  For each answer we
also assess answer quality, groundedness in the visible passages, and whether
the evidence makes the result diagnosable.

\subsection{Failure Categories}
Failures are separated into four types: irrelevant retrieval, insufficient
corpus coverage, ambiguous query, and generation hallucination.  This
distinction prevents attributing an incorrect answer to the generator when
the index never contained suitable evidence.

\section{Evaluation Query Set}
Table~\ref{tab:queries} lists the nine controlled evaluation questions.  The
first three good queries preserve PA2 comparability.  The additional topics
were selected only after they were observed in stored PA2 extracted content
or PA2 retrieval outputs.  The three adverse cases are retained from PA2
because they expose ambiguity and corpus-boundary failures.

\begin{table*}[t]
\centering
\small
\caption{PA3 controlled query set and intended diagnostic role}
\label{tab:queries}
\begin{tabular}{p{2.2cm}p{6.0cm}p{7.0cm}}
\toprule
\textbf{Type} & \textbf{Query} & \textbf{Coverage basis / expected behavior} \\
\midrule
Good & Kateri so simptomi visokega krvnega tlaka in zakaj je merjenje pomembno? &
PA2 reranking surfaced a direct hypertension-management article. \\
Good & Kako zdrav življenjski slog pomaga pri sladkorni bolezni tipa 2? &
PA2 surfaced a diabetes and cardiometabolic lifestyle article. \\
Good & Kako izboljšati spanec in zakaj je reden spanec pomemben? &
PA2 retrieval and validation content contain sleep-focused articles. \\
Good & Kateri so značilni simptomi Parkinsonove bolezni in kako pomaga gibanje? &
The extracted PA2 validation sample contains a Parkinson disease article. \\
Good & Kateri znaki lahko kažejo na pomanjkanje vitamina B12? &
The extracted PA2 validation sample contains a vitamin B12 article. \\
Good & Zakaj so beljakovine pomembne pri ohranjanju mišične mase? &
A saved PA2 retrieval run contains protein and muscle-preservation passages. \\
\midrule
Bad & Boli me &
Underspecified symptom; retrieval cannot establish the user's intended pain. \\
Bad & Kakšne so pravne posledice neplačila prometne kazni v Sloveniji? &
Legal question outside the medical corpus; RAG should decline evidence-based answering. \\
Bad & ATM &
Unspecified acronym; PA2 already showed unrelated retrievals and very low reranker scores. \\
\bottomrule
\end{tabular}
\end{table*}

\section{Evaluation Results}
The final run was executed on May 30, 2026, with the fixed retrieval and
generation configuration described above.  The restored PA2 evaluation
database contained 36,816 long segments with embeddings.  The complete trace,
including all retrieved chunks and full generated answers, is stored in
\texttt{pa3/rag/runs/20260530T022922Z\_evaluation.json}.

Table~\ref{tab:results} summarizes the generated answers and manual scores.
The score order is retrieval relevance / RAG answer quality / LLM-only answer
quality.  In the six good cases, retrieval relevance reached the maximum score
of 12/12.  RAG answer quality reached 10/12, while LLM-only answers reached
6/12.  Across all nine cases, RAG answer quality was 15/18 and LLM-only answer
quality was 7/18.

\begin{table*}[t]
\centering
\scriptsize
\caption{Comparison of RAG and LLM-only answers in the final PA3 evaluation}
\label{tab:results}
\setlength{\tabcolsep}{2pt}
\begin{tabular}{p{1.3cm}p{3.0cm}p{4.4cm}p{3.9cm}p{1.4cm}p{2.3cm}}
\toprule
\textbf{Type} & \textbf{Query} & \textbf{RAG answer summary} &
\textbf{LLM-only answer summary} & \textbf{Score} & \textbf{Comment} \\
\midrule
Good & Blood pressure symptoms and measurement &
Correctly explains that hypertension is often asymptomatic, lists severe
warning symptoms, and ties measurement to prevention of stroke, heart and
kidney complications. &
Mentions asymptomatic pressure and complications, but contains incorrect or
very poor wording. &
2/2/1 & Retrieval adds direct evidence for why measurement matters. \\

Good & Type 2 diabetes and lifestyle &
Uses relevant evidence about diet, physical activity, weight loss, glucose
control and possible remission. &
Gives a short generic answer about diet, activity and insulin sensitivity. &
2/1/1 & Retrieval succeeds; one generated sentence about regular treatment is
poorly grounded. \\

Good & Sleep quality and regular sleep &
Grounds the answer in consistent sleep schedule, circadian rhythm, sleep
quality and health risks of chronic sleep deprivation. &
Lists general sleep-hygiene advice with weak Slovenian phrasing. &
2/2/1 & RAG adds concrete reasons and visible sources. \\

Good & Parkinson symptoms and exercise &
Finds the Parkinson article and explains movement, balance, boxing and social
exercise support. &
Provides only a brief generic description of symptoms and movement therapy. &
2/1/1 & Evidence is relevant, but generated terminology is sometimes awkward. \\

Good & Vitamin B12 deficiency signs &
Uses B12-specific passages about tongue changes, mouth symptoms,
paresthesias and neurological issues. &
Lists common signs, but remains generic and has language errors. &
2/2/1 & RAG improves precision and source traceability. \\

Good & Protein and muscle mass &
Explains amino acids, muscle preservation during weight loss and the need to
combine protein intake with strength training. &
States only that proteins are building blocks for muscle growth and repair. &
2/2/1 & Retrieval adds the key nuance that protein alone is not enough. \\

Bad & Boli me &
Shows several unrelated pain forum cases and cautiously states that diagnosis
is not possible from the evidence. &
Responds empathetically and advises medical help, but gives no explanation. &
1/1/1 & Query is too ambiguous for reliable retrieval. \\

Bad & Legal consequences of unpaid traffic fine &
Correctly refuses to infer a legal answer because retrieved sources do not
cover the question. &
Gives confident legal claims without evidence. &
0/2/0 & Medical corpus has insufficient legal coverage; refusal is correct. \\

Bad & ATM &
Shows unrelated financial/spam forum chunks and refuses to answer the
underspecified acronym. &
Hallucinates an incorrect Slovenian expansion and description. &
0/2/0 & Ambiguous acronym produces irrelevant retrieval; RAG remains
diagnosable. \\
\bottomrule
\end{tabular}
\end{table*}

\section{Discussion and Limitations}
The results show the intended trade-off clearly.  When the query matches the
medical corpus, RAG improves answer quality by adding concrete Slovenian
evidence and making each claim inspectable.  The blood-pressure, sleep,
vitamin B12 and protein queries are the strongest examples: the LLM-only
answers are fluent enough to sound plausible, but they are shorter, less
specific, and sometimes linguistically unreliable.

Retrieval fails in two different ways.  First, a query can be underspecified:
``Boli me'' and ``ATM'' do not provide enough intent, so reranking cannot
recover a stable answer from the corpus.  Second, the corpus can be out of
domain: the traffic-fine query retrieves weak legal/forum fragments from a
medical site, but not enough evidence for a legal answer.  In those cases, the
best RAG behavior is not a better answer, but a transparent refusal.

Generation errors still occur after good retrieval.  In the diabetes and
Parkinson cases, relevant chunks were retrieved, but the model introduced
awkward or weakly grounded wording.  This confirms that retrieval quality
constrains the answer, but does not fully guarantee generation quality.

The primary experimental limitation is qualitative assessment: nine questions
provide interpretable case analysis rather than a statistically representative
benchmark.  This is appropriate to the assignment's focus on transparency,
but it should not be interpreted as clinical performance evaluation.

\section{Conclusion}
PA3 extends the strong PA2 retrieval system without disrupting its documented
choices.  It adds a local Ollama generation layer, an evidence-restricted RAG
prompt, an LLM-only comparison mode, and complete retrieval traces for
evaluation.  The central explainability benefit is concrete: every RAG
answer can be checked against the passages that shaped it, and failures can
be classified as missing evidence, poor retrieval, ambiguous input, or model
generation error.

\begin{thebibliography}{1}
\bibitem{ollamagemma}
Ollama, ``gemma3 model library,'' \url{https://ollama.com/library/gemma3},
accessed May 25, 2026.
\end{thebibliography}

\end{document}
